\scanpage{119}
% line head | text | 68
\mspage{68}
% line 1–2 | text | Donc la donnée d'un triple
Donc la \deleted{\unsure{coup.}}\inserted{donnée d'un} triple $(*)$ équivaut à celle d'un
triple
% line 3–5 | display | (g_U prime,g_T prime,g_X)
\[
\begin{aligned}
(g_{U'},g_{T'},g_X)\ \big|\ &g_X\text{ une flèche }X\to X''\text{ correspondant}\\
&\text{à }g_{\unsure{U|T}}=g_{U'}\circ f_U\text{ et }g_T=g_{T'}\circ f_T\text{ déjà}\\
&\deleted{\text{\unsure{déterminés}}}\ \text{\inserted{\unsure{donné|donnés}} (via }g_{U'},g_{T'}\text{),}
\end{aligned}
\]
% line 6 | text | donc en vertu de 1.3.
donc en vertu de 1.3. cela correspond à
% line 7–10 | display | (*) rewritten triple, with evaluation under the bracket
\[
\begin{aligned}
(g_{U'},g_{T'},u)\quad\left|\begin{aligned}
&g_{U'}:U'\to X''\\
&g_{T'}:T'\unsure{-|\to}X''\\
&u:H\longrightarrow
\underset{\scriptstyle (u,t)\longmapsto\operatorname{Hom}_{X''}(g_{U'}f_U(u),g_{T'}f_T(t))}
{\underline{H_{X''}\bigl((g_{U'}\circ f_U)^{\mathrm o}\times(g_{T'}\circ f_T)}} .
\end{aligned}\right.
\end{aligned}\tag*{(*)}
\]
% line 11–12 | text | On va interpréter la donnée de u
On va interpréter la donnée de $u$, en termes d'un
foncteur
% line 13 | display | (**) F upper star
\[
F^*=(H'\mapsto H'\circ F):
\underline{\operatorname{Hom}}({U'}^{\mathrm o}\times T',\mathrm{Ens})\to
\underline{\operatorname{Hom}}(U^{\mathrm o}\times T,\mathrm{Ens})\tag*{(**)}
\]
% line 14 | text | où F is given by
où $F:U^{\mathrm o}\times T\longrightarrow {U'}^{\mathrm o}\times T'$ est donné par $F=f_U^{\mathrm o}\times f_T$.

% line 15–16 | text | On a un objet H double prime
On a un objet $H''$ \inserted{$=H''_{g_{U'},g_{T'}}$} de la catégorie source de $(**)$
déterminé \deleted{par} en termes du couple $(g_{U'},g_{T'})$, par
% line 17 | display | H double prime evaluated at u prime and t prime
\[
H''(u',t')=\operatorname{Hom}_{X''}(g_{U'}(u'),g_{T'}(t')),
\]
% line 18–19 | text | et le bifoncteur but
et le bifoncteur but de la flèche $u$ dans $(*)$ s'interprète comme
% line 20–21 | display | F upper star of H double prime
\[
\begin{aligned}
F^*(H'')=H''\circ F:(u,t)\unsure{\vdash|\mapsto}{}&
 H''(f_U(u),f_T(t))\deleted{\unsure{=}}\\
&\underset{\operatorname{Hom}(g_{U'}f_U(u),g_{T'}f_T(\unsure{t|t'}))}{\shortparallel}.
\end{aligned}
\]
% line 22 | text | Ainsi, u s'interprète
Ainsi, $u$ s'interprète comme une flèche
% line 23 | display | u to F upper star of H double prime
\[
u:H\longrightarrow F^*(H'').
\]
% line 24–25 | text | Or le foncteur F upper star
Or le foncteur $F^*$ a un adjoint à gauche
$F_!$, \deleted{\unsure{et}}\inserted{i.e.} on a
% line 26 | display | (***) Hom adjunction
\[
\operatorname{Hom}(H,F^*(H''))\simeq\operatorname{Hom}(F_!(H),H'')\tag*{(***)}
\]
% archive sheet number: 119
